\section{Convention crosswalk}
\label{app:crosswalk}

The literature uses closely related minor conventions.  We record only the translations used in this paper.

\begin{lemma}[Operation-by-operation crosswalk]
For surjective morphisms between codes already containing $\varnothing$, the older and mandatory-empty-word definitions agree.  An older morphic image with constant-zero or constant-one output coordinates may first delete those coordinates; after adjoining $\varnothing$, the morphism extends by $f(\varnothing)=\varnothing$.  Replacing a code $\A$ by a trunk $T=\Tk_\A(\alpha)$ corresponds to the mandatory-empty-word morphism
\[
 r_T:\A\cup\{\varnothing\}\longrightarrow T\cup\{\varnothing\},
 \qquad
 r_T(c)=
 \begin{cases}
 c,&c\in T,\\
 \varnothing,&c\notin T.
 \end{cases}
\]
Thus the representable-oriented-matroid characterization may be used under the convention of this paper after removal of constant coordinates.
\end{lemma}

\begin{proof}
For a morphism between mandatory-empty-word codes, surjectivity to $\varnothing$ rules out full inverse images of proper target trunks, and nonemptiness of a target trunk rules out empty inverse images.  Conversely, the empty and full trunks have empty and full inverse images automatically.  Constant coordinates are isomorphism-level redundancies in the older convention.  For trunk replacement, the inverse image of a proper trunk rooted at $\beta$ is the trunk rooted at $\alpha\cup\beta$; the map is surjective and fixes the words of $T$.
\end{proof}
